%% toc: P2 10 10.5
\section{Proofs for connective idioms and first-order expressibility}
\label{sec:annex-p02-w02-p02-005-idioms}

Valuations take values in \(\{0,1\}\).  For the first-order items the
notation is that of \cref{def:p02-w02-p02-005-bw-satisfaction}; the closed
term \(n1\) and the term \(nx\) are the abbreviations fixed after
\cref{def:p02-w02-p02-005-bw-syntax}.  In a structure whose addition is
associative, the value of \(nx\) at \(a\) is the \(n\)-fold sum
\(a+\cdots+a\), written \(na\).  Reduct invariance
(\cref{lem:reduct-invariance}) is used without further mention: a
structure and its expansion or reduct satisfy the same sentences of the
smaller language.  The metatheory is ZFC.

\begin{annexitem}{annex:p02-w02-p02-005-idiom-equivalences}{thm:p02-w02-p02-005-idiom-equivalences}{Correctness of the idiom rows}
\begin{proof}
Fix a valuation \(v\) and compute with the truth conditions of the
connectives.

Clause 1.  \(v(A\to B)=0\) iff \(v(A)=1\) and \(v(B)=0\), iff
\(v(\neg B)=1\) and \(v(\neg A)=0\), iff \(v(\neg B\to\neg A)=0\).  So the
two formulas have the same value under every valuation.  If \(v(A)=0\) and
\(v(B)=1\), then \(v(A\to B)=1\) and \(v(B\to A)=0\), so \(A\to B\) and
\(B\to A\) are not equivalent.

Clause 2.  \(v(\neg\neg B)=1-(1-v(B))=v(B)\), so the antecedents of
\(\neg\neg B\to A\) and \(B\to A\) have the same value and the conditionals
agree.  Next, \(v(B\to A)=0\) iff \(v(B)=1\) and \(v(A)=0\), iff
\(v(\neg B)=0\) and \(v(A)=0\), iff \(v(\neg B\lor A)=0\).

Clause 3.  \(v(\neg A\land\neg B)=1\) iff \(v(A)=0\) and \(v(B)=0\), iff
\(v(A\lor B)=0\), iff \(v(\neg(A\lor B))=1\).

Clause 4.  If \(v(\neg A\land\neg B)=1\), then \(v(A)=0\), so
\(v(A\land B)=0\) and \(v(\neg(A\land B))=1\).  Now let \(v\) make exactly
one of \(A\), \(B\) true.  Then \(v(A\land B)=0\), so
\(v(\neg(A\land B))=1\); and one of \(\neg A\), \(\neg B\) has value \(0\),
so \(v(\neg A\land\neg B)=0\).  Thus \(\neg(A\land B)\) and
\(\neg A\land\neg B\) differ at \(v\).  Finally \(v(\neg A\lor\neg B)=1\)
iff \(v(A)=0\) or \(v(B)=0\), iff \(v(A\land B)=0\); so
\(\neg A\lor\neg B\) has the value of \(\neg(A\land B)\) at every
valuation, and in particular also differs from \(\neg A\land\neg B\) at
\(v\).
\end{proof}
\end{annexitem}

\begin{annexitem}{annex:p02-w02-p02-005-bw-compactness}{thm:p02-w02-p02-005-bw-compactness}{Compactness in two forms}
\begin{proof}
Clause 1 is the right-to-left half of clause 1 of
\cref{thm:fol-compactness}.

Clause 1 implies clause 2; this is the argument of the source
\cite[ch.~A.1, \S2, p.~10]{Barwise1977FirstOrder}.  Suppose
\(T\models\psi\).  Then \(T\cup\{\neg\psi\}\) has no model, since a model
of \(T\) satisfies \(\psi\).  By clause 1 some finite
\(S\subseteq T\cup\{\neg\psi\}\) has no model.  Put \(T_0=S\cap T\), a
finite subset of \(T\) with \(S\subseteq T_0\cup\{\neg\psi\}\).  A model
of \(T_0\cup\{\neg\psi\}\) would be a model of \(S\), so there is none.
Hence every model of \(T_0\) satisfies \(\psi\): \(T_0\models\psi\).

Clause 2 implies clause 1.  Suppose \(T\) has no model, and let \(\psi\)
be \(\exists x\,(x\neq x)\).  Then \(T\models\psi\) holds vacuously, so by
clause 2 some finite \(T_0\subseteq T\) satisfies \(T_0\models\psi\).
No structure satisfies \(\psi\), because \(a=a\) for every element \(a\).
So \(T_0\) has no model.  Contraposing: if every finite subset of \(T\) has
a model, then so does \(T\).

Barwise's route to clause 1, through the Main Lemma
\cite[ch.~A.1, \S4, p.~33]{Barwise1977FirstOrder}.  Let every finite subset
of \(T\) have a model.  By (iii)\(\Rightarrow\)(i) of
\cref{thm:p02-w02-p02-005-bw-main-lemma} it suffices that
\(T\cup T_{\mathrm{Henkin}}\cup\mathrm{Eq}\) be consistent in the sense
of propositional logic, and by \cref{thm:p02-w02-p02-005-bw-prop-compactness}
it suffices that each finite subset be so.  A finite subset lies in
\(T_0\cup T_{\mathrm{Henkin}}\cup\mathrm{Eq}\) for a finite
\(T_0\subseteq T\); \(T_0\) has a model, so this set is consistent by
(i)\(\Rightarrow\)(iii) of the Main Lemma, and so is its subset.
\end{proof}
\end{annexitem}

\begin{annexitem}{annex:p02-w02-p02-005-bw-ls}{thm:p02-w02-p02-005-bw-ls}{L\"owenheim--Skolem}
\begin{proof}
Let \(\mathfrak M\) be a model of \(T\), where \(T\) has at most
\(\kappa\) members.  Let \(L_T\) be the set of symbols of \(L\) occurring
in some member of \(T\).  Each sentence is a finite string, so \(L_T\) is
a union of at most \(\kappa\) finite sets and \(\abs{L_T}\leq\kappa\cdot
\aleph_0=\kappa\).  Let \(\mathfrak M_0\) be the reduct of \(\mathfrak M\)
to \(L_T\); it is a model of \(T\).  Apply \cref{thm:downward-ls} to
\(\mathfrak M_0\) with \(A=\emptyset\).  It gives
\(\mathfrak N_0\preccurlyeq\mathfrak M_0\) whose domain \(N\) satisfies
\(\abs{N}\leq\max(0,\abs{L_T},\aleph_0)\leq\kappa\).  Each member of \(T\)
is a sentence of \(L_T\) true in \(\mathfrak M_0\), hence true in the
elementary substructure \(\mathfrak N_0\).  The domain \(N\) is nonempty;
fix \(n_0\in N\).  Expand \(\mathfrak N_0\) to an \(L\)-structure
\(\mathfrak N\) by interpreting each relation symbol of \(L-L_T\) as the
empty relation, each function symbol of \(L-L_T\) as the constant function
with value \(n_0\), and each constant symbol of \(L-L_T\) as \(n_0\).  Then
\(\mathfrak N\) is a model of \(T\) with domain of cardinality at most
\(\kappa\).

For the last sentence, let \(L\) be finite or countable.  Every sentence of
\(L\) is a finite string over the countable set consisting of \(L\), the
countably many variables and the finitely many other logical symbols, and
there are countably many such strings.  So \(T\) has at most \(\aleph_0\)
members, and the case \(\kappa=\aleph_0\) gives a model of cardinality at
most \(\aleph_0\).

Barwise's route, through the Main Lemma
\cite[ch.~A.1, \S4, p.~33]{Barwise1977FirstOrder}.  Reduce as above to the
language \(L_T\), which has at most \(\kappa\) symbols, and form its
witnessing expansion \(L_T(C)\).  Every formula of \(L_T\) is a finite
string over at most \(\kappa\) symbols, so there are at most \(\kappa\)
formulas and \(\abs{C_1}\leq\kappa\).  If \(\abs{C_n}\leq\kappa\), then
\(L_n\) has at most \(\kappa\) symbols and at most \(\kappa\) formulas, so
\(\abs{C_{n+1}}\leq\kappa\); hence \(\abs C\leq\aleph_0\cdot\kappa=\kappa\).
By (i)\(\Rightarrow\)(ii) of \cref{thm:p02-w02-p02-005-bw-main-lemma},
\(T\) has a canonical model for \(L_T(C)\), whose domain
\(\{c^{\mathfrak M}\mid c\in C\}\) has at most \(\kappa\) elements.  Its
reduct to \(L_T\), expanded to \(L\) as above, is the required model.
\end{proof}
\end{annexitem}

\begin{annexitem}{annex:p02-w02-p02-005-bw-divisible}{thm:p02-w02-p02-005-bw-divisible}{Divisibility is not finitely axiomatizable}
\begin{proof}
The route is the source's \cite[ch.~A.1, \S2, pp.~10--11]{Barwise1977FirstOrder}.
Let \(\{\psi_1,\ldots,\psi_k\}\) be true in all divisible abelian groups,
and let \(\psi\) be \(\psi_1\land\cdots\land\psi_k\), or
\(\forall x\,(x=x)\) when \(k=0\).  Let \(T\) consist of (1)--(4) and all
\((6)_n\) for \(n\geq2\).  The models of \(T\) are exactly the divisible
abelian groups, because \((6)_1\) holds in every structure.  So
\(T\models\psi\), and clause 2 of
\cref{thm:p02-w02-p02-005-bw-compactness} gives a finite
\(T_0\subseteq T\) with \(T_0\models\psi\).  Choose \(N\geq2\) such that
every \((6)_n\) in \(T_0\) has \(n\leq N\).  Every abelian group satisfying
\((6)_n\) for \(2\leq n\leq N\) is then a model of \(T_0\), hence of
\(\psi\).

Choose a prime \(p>N\) and let \(\mathbb Z_p\) be the integers modulo
\(p\) under addition modulo \(p\), with identity \(0\).  It is an abelian
group.  For \(2\leq n<p\), the map \(a\mapsto na\) on \(\mathbb Z_p\) is
injective: if \(na=nb\), then \(p\) divides \(n(a-b)\); \(p\) is prime and
does not divide \(n\), so \(p\) divides \(a-b\) and \(a=b\).  An injective
map of a finite set to itself is surjective, so for every \(a\) there is
\(b\) with \(nb=a\), and \(\mathbb Z_p\models(6)_n\).  Since \(N<p\),
\(\mathbb Z_p\) satisfies \((6)_n\) for \(2\leq n\leq N\), so
\(\mathbb Z_p\models\psi\), that is, \(\mathbb Z_p\) satisfies each
\(\psi_i\).  But \(pb=0\) for every \(b\in\mathbb Z_p\), and \(1\neq0\)
because \(p\geq2\); so no \(b\) satisfies \(pb=1\), \((6)_p\) fails, and
\(\mathbb Z_p\) is not divisible.
\end{proof}
\end{annexitem}

\begin{annexitem}{annex:p02-w02-p02-005-bw-torsion-free}{thm:p02-w02-p02-005-bw-torsion-free}{Torsion-freeness is not finitely axiomatizable}
\begin{proof}
The source states that the proof has the form of the previous one and does
not write it \cite[ch.~A.1, \S2, p.~11]{Barwise1977FirstOrder}; here it is
in full.  Suppose a finite set \(\Phi\) of sentences of \(\{+,0\}\) has as
its models exactly the torsion-free abelian groups.  Let \(\psi\) be the
conjunction of \(\Phi\), or \(\forall x\,(x=x)\) if \(\Phi\) is empty.
Let \(T\) consist of (1)--(4) and all \((7)_n\) for \(n\geq1\).  The
models of \(T\) are exactly the torsion-free abelian groups, so they are
exactly the models of \(\Phi\), and \(T\models\psi\).  Clause 2 of
\cref{thm:p02-w02-p02-005-bw-compactness} gives a finite
\(T_0\subseteq T\) with \(T_0\models\psi\).  Choose \(N\geq1\) bounding
every \(n\) with \((7)_n\in T_0\), and a prime \(p>N\).

In \(\mathbb Z_p\), let \(a\neq0\) and \(1\leq n<p\).  If \(na=0\), then
\(p\) divides \(na\); since \(p\) is prime and divides neither \(n\) nor
\(a\) (as \(0<a<p\)), this is impossible.  So \(\mathbb Z_p\models(7)_n\)
for \(n\leq N\), hence \(\mathbb Z_p\models T_0\), hence
\(\mathbb Z_p\models\psi\) and \(\mathbb Z_p\) is a model of \(\Phi\).
Then \(\mathbb Z_p\) is torsion-free.  But \(1\neq0\) and \(p1=0\) in
\(\mathbb Z_p\), so \((7)_p\) fails.  This contradiction shows that no such
\(\Phi\) exists.
\end{proof}
\end{annexitem}

\begin{annexitem}{annex:p02-w02-p02-005-bw-torsion}{thm:p02-w02-p02-005-bw-torsion}{Torsion is not axiomatizable}
\begin{proof}
The route is the source's \cite[ch.~A.1, \S2, p.~11]{Barwise1977FirstOrder}.
Here an element \(a\) of an abelian group has \emph{order \(\geq n\)} if
\(ma\neq0\) for \(1\leq m<n\).

\emph{The second sentence.}  Let \(G=(G,+,0)\) be abelian and, for each
\(n\), let \(x_n\in G\) have order \(\geq n\).  Add a new constant symbol
\(c\).  Let \(T\) consist of every sentence of \(\{+,0\}\) true in \(G\)
together with the sentences \(nc\neq0\) for \(n\geq2\).  Let
\(T_0\subseteq T\) be finite, and choose \(N\geq2\) larger than every
\(n\) with \(nc\neq0\in T_0\).  Interpret \(c\) as \(x_N\).  Since \(x_N\)
has order \(\geq N\), \(nx_N\neq0\) for \(2\leq n<N\), and
\((G,+,0,x_N)\) is a model of \(T_0\).  By clause 1 of
\cref{thm:p02-w02-p02-005-bw-compactness}, \(T\) has a model
\((H,+,0,c^H)\).

Let \(H=(H,+,0)\) be its reduct.  For each sentence \(\sigma\) of
\(\{+,0\}\), exactly one of \(\sigma\), \(\neg\sigma\) is true in \(G\) and
so belongs to \(T\); hence \(H\models\sigma\) iff \(G\models\sigma\), and
\(G\equiv H\).  In particular \(H\) satisfies (1)--(4) and is an abelian
group.  It is not torsion.  Suppose \(mc^H=0\) for some \(m\geq1\).  If
\(m\geq2\), this contradicts \(mc\neq0\in T\).  If \(m=1\), then
\(c^H=0\) and \(2c^H=0+0=0\), contradicting \(2c\neq0\in T\).

\emph{The first sentence.}  Let \(G\) be the direct sum of the groups
\(\mathbb Z_p\) over all primes \(p\): the functions \(g\) on the primes
with \(g(p)\in\mathbb Z_p\) and \(g(p)=0\) for all but finitely many
\(p\), added coordinatewise.  It is an abelian group.  It is torsion: if
\(g(p)=0\) outside the finite set \(F\) of primes, then \(mg=0\) for \(m\)
the product of the members of \(F\) (and \(m=1\) when \(F\) is empty).
For a prime \(p\), the element \(e_p\) with \(e_p(p)=1\) and value \(0\)
elsewhere satisfies \(me_p\neq0\) for \(1\leq m<p\).  Given \(n\), choose a
prime \(p\geq n\); then \(e_p\) has order \(\geq n\).  The second sentence
gives an abelian group \(H\equiv G\) that is not torsion.  Every sentence
true in all torsion abelian groups is true in \(G\), hence in \(H\).

\emph{The last sentence.}  If a set \(\Phi\) of sentences had exactly the
torsion abelian groups as models, then \(G\models\Phi\), so
\(H\models\Phi\) because \(H\equiv G\), and \(H\) would be torsion.
\end{proof}
\end{annexitem}

\begin{annexitem}{annex:p02-w02-p02-005-bw-reals}{thm:p02-w02-p02-005-bw-reals}{No first-order categorical axioms for the reals}
\begin{proof}
The route is the source's \cite[ch.~A.1, \S2, p.~12]{Barwise1977FirstOrder},
with the cardinality comparison written out.  The language
\(\{+,\cdot,<,0,1\}\) is finite.  The set of sentences true in
\(\mathbb R\) has the model \(\mathbb R\), so by
\cref{thm:p02-w02-p02-005-bw-ls} it has a model \(\mathfrak A\) whose
domain \(A\) has at most \(\aleph_0\) elements.

The set \(\mathbb R\) is not countable.  Let \(r_0,r_1,\ldots\) be any
sequence of reals.  Choose closed intervals
\(I_0\supseteq I_1\supseteq\cdots\) of positive length with
\(r_k\notin I_k\): given \(I_{k-1}=[a,b]\) (or \(I_{-1}=[0,1]\)), split it
into its closed thirds; the point \(r_k\) lies in at most two of the
three, whose union is not the whole interval, so some third omits
\(r_k\).  By the nested interval property of \(\mathbb R\), some real lies
in every \(I_k\), and it differs from every \(r_k\).  So no sequence
enumerates \(\mathbb R\).

Now suppose \(\Phi\) is a set of sentences whose models are exactly the
structures isomorphic to \(\mathbb R\).  Then \(\mathbb R\models\Phi\), so
every member of \(\Phi\) is true in \(\mathbb R\) and therefore in
\(\mathfrak A\).  So \(\mathfrak A\) is isomorphic to \(\mathbb R\), and
\(A\) is in bijection with \(\mathbb R\).  This contradicts
\(\abs{A}\leq\aleph_0\).
\end{proof}
\end{annexitem}

\begin{annexitem}{annex:p02-w02-p02-005-bw-nonstandard}{thm:p02-w02-p02-005-bw-nonstandard}{A non-Archimedean elementary extension}
\begin{proof}
The route is the source's \cite[ch.~A.1, \S2, p.~12]{Barwise1977FirstOrder};
the embedding of \(\mathbb R\) is written out.  Let \(L_{\mathbb R}\) be
\(\{+,\cdot,<,0,1\}\) together with a constant symbol \(\underline r\) for
each real \(r\), and let \(\mathbb R_{\mathbb R}\) be the expansion of
\(\mathbb R\) with \(\underline r\) interpreted as \(r\).  (The source
uses \(r\) as its own name; the underline only separates symbol from
number.)  Add a new constant symbol \(c\), and let \(T\) consist of every
sentence of \(L_{\mathbb R}\) true in \(\mathbb R_{\mathbb R}\) together
with \(\underline r<c\) for every real \(r\).

A finite \(T_0\subseteq T\) contains \(\underline r<c\) for only finitely
many \(r\), say \(r_1,\ldots,r_m\).  Interpret \(c\) in
\(\mathbb R_{\mathbb R}\) as \(1+\max(r_1,\ldots,r_m)\), or as \(0\) if
\(m=0\); this models \(T_0\).  By clause 1 of
\cref{thm:p02-w02-p02-005-bw-compactness}, \(T\) has a model
\(\mathfrak A\).

Define \(h:\mathbb R\to A\) by \(h(r)=\underline r^{\mathfrak A}\).  For
reals \(r,s,t\): if \(r\neq s\), then \(\underline r\neq\underline s\) is
in \(T\); if \(r+s=t\), then \(\underline r+\underline s=\underline t\) is
in \(T\); if \(r\cdot s=t\), then \(\underline r\cdot\underline s=
\underline t\) is in \(T\); \(\underline r<\underline s\) is in \(T\) if
\(r<s\), and its negation is in \(T\) otherwise; and
\(\underline0=0\), \(\underline1=1\) are in \(T\).  So \(h\) is injective
and preserves \(+\), \(\cdot\), \(0\), \(1\), and \(<\) in both directions:
it embeds \(\mathbb R\) into the \(\{+,\cdot,<,0,1\}\)-reduct of
\(\mathfrak A\).  Replace each \(h(r)\) by \(r\) (after renaming, if
necessary, the elements of \(A\) outside the image so that they are not
reals) and call the resulting reduct \({}^*\mathbb R\).  It contains
\(\mathbb R\) as a substructure, and it satisfies every sentence of
\(L_{\mathbb R}\) true in \(\mathbb R_{\mathbb R}\) when \(\underline r\)
is read as \(r\).

The ordered-field axioms are finitely many first-order sentences true in
\(\mathbb R\), so \({}^*\mathbb R\) is an ordered field.  For each
\(n\geq1\), the sentence \(n1=\underline n\) is true in
\(\mathbb R_{\mathbb R}\), so \(n1=n\) in \({}^*\mathbb R\); and
\(\underline n<c\) is in \(T\), so \(n<c^{\mathfrak A}\).  Thus
\(c^{\mathfrak A}\) exceeds every \(n1\), and \({}^*\mathbb R\) is not
Archimedean.
\end{proof}
\end{annexitem}

\begin{annexitem}{annex:p02-w02-p02-005-bw-maximal}{thm:p02-w02-p02-005-bw-maximal}{A first-order equivalent of maximality}
\begin{proof}
Here an \emph{ideal} in the broad sense is an additive subgroup \(J\) with
\(x\cdot y\in J\) whenever \(y\in J\); \(I\) is maximal when every ideal
\(J\supseteq I\) in this sense equals \(I\) or \(R\).  The source obtains
(9)\('\) by saying that \(R/I\) is a field, using the ring-theoretic lemma
that \(I\) is maximal iff \(R/I\) is a field
\cite[ch.~A.1, \S2, p.~14]{Barwise1977FirstOrder}: (9)\('\) states that
each coset \(x+I\neq0+I\) has an inverse \(y+I\).  The following proof is
direct and does not presuppose that lemma.

Suppose \((\mathfrak R,I)\models(9)'\).  Let \(J\supseteq I\) be an
ideal with \(J\neq I\), and pick \(x\in J-I\).  By (9)\('\) there are
\(y\in R\) and \(z\in I\) with \(xy+z=1\).  Then \(xy=yx\in J\) by
absorption and \(z\in J\), so \(1\in J\), and \(r=r\cdot1\in J\) for
every \(r\in R\).  So \(J=R\), and \(I\) is maximal.

Conversely let \(I\) be maximal and \(x\notin I\).  Put
\(J=\{xy+z:y\in R,\ z\in I\}\).  It contains \(0=x0+0\).  It is closed
under addition, since \((xy+z)+(xy'+z')=x(y+y')+(z+z')\), and under
negation, since \(-(xy+z)=x(-y)+(-z)\) with \(-z\in I\) by the third ideal
axiom.  It absorbs multiplication, since
\(r(xy+z)=x(ry)+rz\) by commutativity and associativity, with
\(rz\in I\).  It contains \(I\) (take \(y=0\)) and \(x\) (take \(y=1\),
\(z=0\)), so \(J\neq I\).  By maximality \(J=R\), so \(1\in J\): there are
\(y\) and \(z\in I\) with \(xy+z=1\).  As \(x\notin I\) was arbitrary,
\((\mathfrak R,I)\models(9)'\).
\end{proof}
\end{annexitem}

\begin{annexitem}{annex:p02-w02-p02-005-bw-pir}{thm:p02-w02-p02-005-bw-pir}{Being a principal ideal ring is not first-order}
\begin{proof}
The source asserts that a compactness argument gives such a ring and leaves
the argument as an exercise \cite[ch.~A.1, \S2, p.~14]{Barwise1977FirstOrder}.
The following is one such argument.  Write \(2\) for the closed term
\(1+1\).

Add a new constant symbol \(c\).  Let \(T\) consist of all sentences of
\(\{+,\cdot,0,1\}\) true in \(\mathbb Z\), the sentence \(c\neq0\), and the
sentences \(\delta_n\): \(\exists y\,((n1)\cdot y=c)\) for \(n\geq1\).  For
a finite \(T_0\subseteq T\), let \(N\) bound every \(n\) with
\(\delta_n\in T_0\) and interpret \(c\) in \(\mathbb Z\) as \(N!\), which
is nonzero and divisible by each \(n\leq N\); this models \(T_0\).  By
clause 1 of \cref{thm:p02-w02-p02-005-bw-compactness}, \(T\) has a model;
let \(\mathfrak R=(R,+,\cdot,0,1)\) be its reduct and \(c^*\) the value of
\(c\).  As in the proof of \cref{thm:p02-w02-p02-005-bw-torsion},
\(\mathfrak R\equiv\mathbb Z\).  So \(\mathfrak R\) is a commutative ring
with identity, and the following sentences, true in \(\mathbb Z\), hold in
\(\mathfrak R\):
\[
\begin{aligned}
&\neg\exists y\,(2\cdot y=1),\\
&\forall u\,\forall w\,(2\cdot u=2\cdot w\to u=w),\\
&\forall d\,\forall z\,(d\neq0\to d\neq(2\cdot z)\cdot d),\\
&(2n)1=2\cdot(n1)\quad\text{for each }n\geq1.
\end{aligned}
\]
The third holds in \(\mathbb Z\) because \(d=2zd\) with \(d\neq0\) would
give \(2z=1\).

Let \(I\) be the set of \(a\in R\) such that for every \(n\geq1\) some
\(b\in R\) has \((n1)\cdot b=a\).  It contains \(0\) (take \(b=0\)) and
\(c^*\) (by the sentences \(\delta_n\)).  If \((n1)b=a\) and
\((n1)b'=a'\), then \((n1)(b+b')=a+a'\), \((n1)(-b)=-a\) and
\((n1)(rb)=ra\) for every \(r\in R\), by the ring axioms; so \(I\) is
closed under addition and negation and absorbs multiplication.  It omits
\(1\), since \(2\cdot y=1\) has no solution.  So \(I\) is a proper ideal,
and it contains \(c^*\neq0\).

Suppose \(I=\{z\cdot d:z\in R\}\) for some \(d\).  Then \(d=1\cdot d\in
I\), and \(d\neq0\) because \(c^*\in I\) is a multiple of \(d\) and is
nonzero.  As \(d\in I\), \(d=2\cdot e\) for some \(e\).  This \(e\) lies
in \(I\): given \(n\geq1\), \(d\in I\) gives \(k\) with
\(d=((2n)1)\cdot k=2\cdot((n1)\cdot k)\), so \(2\cdot e=2\cdot((n1)\cdot
k)\) and cancellation gives \(e=(n1)\cdot k\).  Hence \(e=z\cdot d\) for
some \(z\), and \(d=2\cdot(z\cdot d)=(2\cdot z)\cdot d\) with \(d\neq0\),
contradicting the third displayed sentence.  So \(I\) is not principal.

For the consequence: if a set \(\Phi\) of sentences had exactly the
principal ideal rings as models, then \(\mathbb Z\models\Phi\), since
every ideal of \(\mathbb Z\) is generated by its least positive element (or
is \(\{0\}\)); so \(\mathfrak R\models\Phi\) because
\(\mathfrak R\equiv\mathbb Z\), and \(\mathfrak R\) would be a principal
ideal ring.
\end{proof}
\end{annexitem}

\begin{annexitem}{annex:p02-w02-p02-005-bw-characteristic}{thm:p02-w02-p02-005-bw-characteristic}{Transfer from characteristic \(0\)}
\begin{proof}
The source says only that the proof is like that of
\cref{thm:p02-w02-p02-005-bw-divisible}
\cite[ch.~A.1, \S2, p.~15]{Barwise1977FirstOrder}; here it is in full.
Let \(F\) consist of the commutative-ring axioms, the field axiom
\(\forall x\,\exists y\,[x\neq0\to x\cdot y=1]\), and \((10)_p\) for every
prime \(p\).  Its models are exactly the fields of characteristic \(0\), so
\(F\models\psi\).  Clause 2 of \cref{thm:p02-w02-p02-005-bw-compactness}
gives a finite \(F_0\subseteq F\) with \(F_0\models\psi\).  Let
\(N_\psi\) be the largest prime \(p\) with \((10)_p\in F_0\), or \(1\) if
there is none.

Let \(K\) be a field of characteristic \(q\), a prime with \(q>N_\psi\);
so \(q1=0\) in \(K\).  Let \(p\neq q\) be prime.  As \(p\) and \(q\) are
distinct primes, there are integers \(a,b\) with \(ap+bq=1\).  In any ring,
integer multiples of \(1\) satisfy \((ap+bq)1=a(p1)+b(q1)\).  If \(p1=0\)
in \(K\), this gives \(1=0\), which the axiom \(0\neq1\) excludes.  So
\(K\models(10)_p\) for every prime \(p\neq q\), in particular for every
\(p\leq N_\psi\).  Hence \(K\models F_0\) and \(K\models\psi\).
\end{proof}
\end{annexitem}

\subsection{Proofs for Barwise's \S\S3--4}
\label{sec:annex-p02-w02-p02-005-bw-s4}

Throughout, \(\varphi(t)\) is \(\varphi(t/v)\) with no variable of \(t\)
bound in \(\varphi\); under this proviso the substitution is the
capture-avoiding substitution of \cref{thm:fol-substitution}, which
therefore applies.  Formulas are uniquely readable, so every first-order
formula is a propositional formula over its prime constituents in exactly
one way.  For a truth assignment \(v\) and a set \(\Theta\) of formulas,
\(v\) \emph{satisfies} \(\Theta\) if \(\bar v(\theta)=\mathbf t\) for all
\(\theta\in\Theta\).

\begin{annexitem}{annex:p02-w02-p02-005-bw-iso-equiv}{thm:p02-w02-p02-005-bw-iso-equiv}{Isomorphic structures are elementarily equivalent}
\begin{proof}
For terms, \(h(t^{\mathfrak M}(s))=t^{\mathfrak N}(h\circ s)\) by induction
on \(t\): for a constant, \(h(c^{\mathfrak M})=c^{\mathfrak N}\); for a
variable, both sides are \(h(s(v))\); for \(f(t_1\cdots t_n)\),
\(h(f^{\mathfrak M}(\ldots))=f^{\mathfrak N}(h(\ldots))\) and the induction
hypothesis.  For formulas induct on \(\varphi\).  Equations: \(h\) is
injective, so \(t_1^{\mathfrak M}(s)=t_2^{\mathfrak M}(s)\) iff their
\(h\)-images agree.  \(R(t_1\cdots t_n)\): \(h\) maps \(R^{\mathfrak M}\)
onto \(R^{\mathfrak N}\) and its complement onto the complement.
Connectives follow from the induction hypothesis.
\(\exists v\,\varphi\): \(\mathfrak M\models\varphi[s(\tfrac{a}{v})]\) for
some \(a\in M\) iff \(\mathfrak N\models\varphi[(h\circ s)(\tfrac{h(a)}{v})]\)
for some \(a\in M\), iff \(\mathfrak N\models\varphi[(h\circ s)(\tfrac{b}{v})]\)
for some \(b\in N\), because \(h\) is onto and
\(h\circ(s(\tfrac av))=(h\circ s)(\tfrac{h(a)}v)\).  \(\forall v\,\varphi\)
is the same with ``all''.  For a sentence the assignment is irrelevant
(\cref{lem:fol-coincidence}), so \(\mathfrak M\equiv\mathfrak N\).
\end{proof}
\end{annexitem}

\begin{annexitem}{annex:p02-w02-p02-005-bw-prop-compactness}{thm:p02-w02-p02-005-bw-prop-compactness}{Compactness for propositional logic}
\begin{proof}
If \(T\) is consistent, a truth assignment satisfying \(T\) satisfies each
finite subset.  For the converse call a set \emph{finitely consistent} if
each of its finite subsets is consistent.

\emph{First proof} \cite[ch.~A.1, \S4, pp.~26--27]{Barwise1977FirstOrder}.
Call \(S\) \emph{maximal finitely consistent} if it is finitely consistent
and contains \(A\) or \(\neg A\) for every formula \(A\).

Step 1: a maximal finitely consistent \(S\) is satisfied by the assignment
\(v_S\) with \(v_S(p)=\mathbf t\) iff \(p\in S\).  First,
(a) \(\neg B\in S\) iff \(B\notin S\): both in \(S\) would make the finite
set \(\{B,\neg B\}\) inconsistent, and neither contradicts maximality.
(b) \((A\land B)\in S\) iff \(A\in S\) and \(B\in S\): if
\((A\land B)\in S\) and \(A\notin S\), then \(\neg A\in S\) and
\(\{A\land B,\neg A\}\) is inconsistent (similarly for \(B\)); if
\(A,B\in S\) and \((A\land B)\notin S\), then \(\{A,B,\neg(A\land B)\}\subseteq S\)
is inconsistent.  (c) \((A\lor B)\in S\) iff \(A\in S\) or \(B\in S\): if
\((A\lor B)\in S\) and neither disjunct, then \(\neg A,\neg B\in S\) and
\(\{A\lor B,\neg A,\neg B\}\) is inconsistent (Barwise's sample case); if
\(A\in S\) and \((A\lor B)\notin S\), then \(\{A,\neg(A\lor B)\}\subseteq S\)
is inconsistent, and likewise for \(B\).  (d) \((A\to B)\in S\) iff
\(A\notin S\) or \(B\in S\): if \((A\to B),A\in S\) and \(B\notin S\), then
\(\{A\to B,A,\neg B\}\subseteq S\) is inconsistent; if \((A\to B)\notin S\),
then \(\neg(A\to B)\in S\), and \(A\notin S\) gives the inconsistent
\(\{\neg A,\neg(A\to B)\}\subseteq S\), while \(B\in S\) gives the
inconsistent \(\{B,\neg(A\to B)\}\).  By induction on \(A\), using
(a)--(d) and the clauses of \(\bar v_S\),
\(\bar v_S(A)=\mathbf t\) iff \(A\in S\).  So \(S=\{A\mid\bar v_S(A)=\mathbf t\}\);
conversely, for every \(v\) the set \(\{A\mid\bar v(A)=\mathbf t\}\) is
maximal finitely consistent.  Hence it suffices to extend a finitely
consistent \(T\) to a maximal finitely consistent \(S\).

Step 2: the extension.  Suppose the formulas are listed as
\(A_\alpha\), \(\alpha<\lambda\), along a well-ordering.  When \(P\) is
countable, \(\lambda=\omega\) and the list is an enumeration; when \(P\) is
well-ordered, the formulas, being finite strings over \(P\) and finitely
many further symbols, are well-ordered by length and then
lexicographically, without the axiom of choice; for arbitrary \(P\) the
axiom of choice gives a well-ordering of \(P\).  Put \(T_0=T\);
\(T_{\alpha+1}=T_\alpha\cup\{A_\alpha\}\) if this is finitely consistent,
and \(T_{\alpha+1}=T_\alpha\cup\{\neg A_\alpha\}\) otherwise;
\(T_\delta=\bigcup_{\alpha<\delta}T_\alpha\) at limits; and
\(S=\bigcup_{\alpha<\lambda}T_\alpha\).  Then \(T\subseteq S\) and each
\(A_\alpha\) or its negation is in \(S\).  Each \(T_\alpha\) is finitely
consistent, by induction.  \(T_0=T\) is the hypothesis.  At a limit, a
finite subset of \(T_\delta\) lies in some \(T_\alpha\), \(\alpha<\delta\),
because the \(T_\alpha\) increase.  At a successor, case 1 is immediate.
In case 2 there is a finite \(T'\subseteq T_\alpha\) with
\(T'\cup\{A_\alpha\}\) inconsistent; if \(T_{\alpha+1}\) were not finitely
consistent there would be a finite \(T''\subseteq T_\alpha\) with
\(T''\cup\{\neg A_\alpha\}\) inconsistent.  But \(T'\cup T''\) is a finite
subset of \(T_\alpha\), so some \(v\) satisfies it; \(v\) makes
\(A_\alpha\) or \(\neg A_\alpha\) true, so it satisfies
\(T'\cup\{A_\alpha\}\) or \(T''\cup\{\neg A_\alpha\}\), a contradiction.
Finally \(S\) is finitely consistent, as a finite subset of \(S\) lies in
some \(T_\alpha\) (for \(\lambda\) a limit; if \(\lambda\) is a successor,
\(S\) is the last \(T_\alpha\)).

\emph{Second proof} \cite[ch.~A.1, \S4, pp.~27--28]{Barwise1977FirstOrder}.
Give \(2=\{\mathbf t,\mathbf f\}\) the discrete topology and
\(X=2^P\), the set of truth assignments, the product topology.  By the
Tychonoff theorem \(X\) is compact (and Hausdorff).  For a formula \(A\) put
\(F_A=\{v\in X\mid\bar v(A)=\mathbf t\}\).  Each \(F_A\) is clopen: for
prime \(p\), \(F_p\) and \(X-F_p=\{v\mid v(p)=\mathbf f\}\) are basic open
sets; and \(F_{\neg A}=X-F_A\), \(F_{A\land B}=F_A\cap F_B\),
\(F_{A\lor B}=F_A\cup F_B\), \(F_{A\to B}=(X-F_A)\cup F_B\), and clopen sets
are closed under complement, finite intersection and finite union.  If
every finite \(T_0\subseteq T\) is consistent, then
\(\bigcap_{A\in T_0}F_A\neq\emptyset\) for each such \(T_0\).  If
\(\bigcap_{A\in T}F_A\) were empty, the open sets \(X-F_A\), \(A\in T\),
would cover \(X\), and a finite subcover would give a finite \(T_0\) with
\(\bigcap_{A\in T_0}F_A=\emptyset\).  So some \(v\) lies in every \(F_A\),
\(A\in T\): it satisfies \(T\).
\end{proof}
\end{annexitem}

\begin{annexitem}{annex:p02-w02-p02-005-bw-map-colouring}{thm:p02-w02-p02-005-bw-map-colouring}{Map colouring and Tychonoff for \(2^P\)}
\begin{proof}
Part 1, along the lines Barwise sketches \cite[ch.~A.1, \S4, p.~28]{Barwise1977FirstOrder}.
Take prime formulas \(p_{x,i}\) for \(x\in X\) and \(1\leq i\leq k\),
read ``country \(x\) has colour \(i\)''.  Let \(T\) consist of
\(p_{x,1}\lor\cdots\lor p_{x,k}\) and \(\neg(p_{x,i}\land p_{x,j})\)
(\(i\neq j\)) for each \(x\in X\), and \(\neg(p_{x,i}\land p_{y,i})\) for
each \((x,y)\in E\) and each \(i\).  A finite \(T_0\subseteq T\) mentions a
finite set \(X_0\) of countries.  Colour the submap on \(X_0\) by \(c\) and
put \(v(p_{x,i})=\mathbf t\) iff \(x\in X_0\) and \(c(x)=i\).  Every member
of \(T_0\) concerns countries of \(X_0\), and adjacent ones among them when
it is an adjacency axiom, so \(v\) satisfies \(T_0\).  By
\cref{thm:p02-w02-p02-005-bw-prop-compactness} some \(v\) satisfies \(T\).
For each \(x\) exactly one \(i\) has \(v(p_{x,i})=\mathbf t\); call it
\(c(x)\).  If \((x,y)\in E\), the axiom
\(\neg(p_{x,c(x)}\land p_{y,c(x)})\) gives \(c(y)\neq c(x)\).

Part 2.  Basic open sets of \(X=2^P\) have the form
\(\{v\mid v(p_1)=e_1,\ldots,v(p_n)=e_n\}=F_B\), where \(B\) is the conjunction
of the literals \(p_j\) (if \(e_j=\mathbf t\)) and \(\neg p_j\) (if
\(e_j=\mathbf f\)); write \(\mathcal B\) for the set of such \(B\).  Let
\(\mathcal C\) be a family of closed sets every finite subfamily of which
has nonempty intersection.  For \(C\in\mathcal C\) put
\(U_C=\{B\in\mathcal B\mid F_B\subseteq X-C\}\); since \(X-C\) is open,
\(X-C=\bigcup_{B\in U_C}F_B\) and \(C=\bigcap_{B\in U_C}F_{\neg B}\).  Let
\(T=\{\neg B\mid B\in U_C,\ C\in\mathcal C\}\).  A finite \(T_0\subseteq T\)
draws on finitely many \(C_1,\ldots,C_m\in\mathcal C\), and any \(v\) in
\(C_1\cap\cdots\cap C_m\neq\emptyset\) satisfies \(T_0\).  By the first
proof of \cref{thm:p02-w02-p02-005-bw-prop-compactness}, which uses no
topology, some \(v\) satisfies \(T\); then \(v\in\bigcap\mathcal C\).  So
\(X\) is compact.  No choice is used beyond that in the first proof.
\end{proof}
\end{annexitem}

\begin{annexitem}{annex:p02-w02-p02-005-bw-lemma-43}{thm:p02-w02-p02-005-bw-lemma-43}{Satisfaction induces a truth assignment}
\begin{proof}
For prime \(\varphi\) put \(v(\varphi)=\mathbf t\) if
\(\mathfrak M\models\varphi[s]\) and \(\mathbf f\) otherwise
\cite[ch.~A.1, \S4, p.~29]{Barwise1977FirstOrder}.  Every formula is built
from prime formulas by \(\neg,\land,\lor,\to\), and clauses (iii)--(vi) of
\cref{def:p02-w02-p02-005-bw-satisfaction} have the same form as the
clauses of \(\bar v\); induction on the propositional construction gives
\(\mathfrak M\models\varphi[s]\) iff \(\bar v(\varphi)=\mathbf t\).  If every
sentence of \(\Sigma\) is true in \(\mathfrak M\), this \(v\) (for any
\(s\)) satisfies \(\Sigma\).  A tautology \(\tau\) has
\(\bar v(\tau)=\mathbf t\) for this \(v\), so \(\mathfrak M\models\tau[s]\).
\end{proof}
\end{annexitem}

\begin{annexitem}{annex:p02-w02-p02-005-bw-43-converse}{thm:p02-w02-p02-005-bw-43-converse}{The converse of Lemma~4.3 fails}
\begin{proof}
The three sentences are distinct prime formulas, so the truth assignment
giving each of them \(\mathbf t\) satisfies the set.  Suppose
\(\mathfrak M\) were a model.  By the third sentence some \(a\in M\) has
\(a\notin S^{\mathfrak M}\); by the second \(a\in R^{\mathfrak M}\); by
the first \(a\in S^{\mathfrak M}\), a contradiction.
\end{proof}
\end{annexitem}

\begin{annexitem}{annex:p02-w02-p02-005-bw-lemma-47}{thm:p02-w02-p02-005-bw-lemma-47}{Every structure expands to a model of \(T_{\mathrm{Henkin}}\)}
\begin{proof}
Barwise's route \cite[ch.~A.1, \S4, pp.~30--31]{Barwise1977FirstOrder}.
Fix a choice function \(g\) on the nonempty subsets of \(M\).  Interpret
the constants of \(C_{n+1}-C_n\) by induction on \(n\).  Suppose the
constants of \(C_n\) are interpreted, giving an expansion \(\mathfrak M_n\)
of \(\mathfrak M\) to \(L_n\).  For \(c=c_{\varphi(v)}\in C_{n+1}-C_n\),
\(\varphi(v)\) is a formula of \(L_n\); put
\(W=\{a\in M\mid\mathfrak M_n\models\varphi[a]\}\) and let
\(c^{\mathfrak M}=g(W)\) if \(W\neq\emptyset\), \(g(M)\) otherwise.  Let
\(\mathfrak M_C\) be the resulting expansion to \(L(C)\); by reduct
invariance it agrees with each \(\mathfrak M_n\) on sentences of \(L_n\).

Type I: if \(\mathfrak M_C\models\exists v\,\varphi(v)\) with
\(c=c_{\varphi(v)}\in C_{n+1}-C_n\), then \(W\neq\emptyset\), so
\(\mathfrak M_n\models\varphi[c^{\mathfrak M}]\), and by
\cref{thm:fol-substitution} and reduct invariance
\(\mathfrak M_C\models\varphi(c)\).  When \(v\notin\FV(\varphi)\), \(W\) is
\(M\) or \(\emptyset\) according as \(\varphi\) is true or false, and the
axiom \(\exists v\,\varphi\to\varphi\) holds because \(M\neq\emptyset\).
Type II: suppose \(\mathfrak M_C\models\varphi(c_{\neg\varphi(v)})\) but
\(\mathfrak M_C\not\models\forall v\,\varphi(v)\).  Then
\(\mathfrak M_C\models\exists v\,\neg\varphi(v)\), and the type I axiom for
\(\neg\varphi(v)\) gives \(\mathfrak M_C\models\neg\varphi(c_{\neg\varphi(v)})\),
a contradiction.  Types III and IV hold in every structure: if
\(\mathfrak M_C\models\forall v\,\varphi(v)\), then
\(\mathfrak M_C\models\varphi[t^{\mathfrak M_C}]\), that is,
\(\mathfrak M_C\models\varphi(t)\) by \cref{thm:fol-substitution}; and if
\(\mathfrak M_C\models\varphi(t)\), then \(t^{\mathfrak M_C}\) witnesses
\(\exists v\,\varphi(v)\).
\end{proof}
\end{annexitem}

\begin{annexitem}{annex:p02-w02-p02-005-bw-main-lemma}{thm:p02-w02-p02-005-bw-main-lemma}{Main Lemma}
\begin{proof}
(ii)\(\Rightarrow\)(i): the reduct to \(L\) of a canonical model of \(T\) is
a model of \(T\) (reduct invariance).

(i)\(\Rightarrow\)(iii): expand a model of \(T\) by
\cref{thm:p02-w02-p02-005-bw-lemma-47} to a model \(\mathfrak M_C\) of
\(T_{\mathrm{Henkin}}\); it is a model of \(T\) by reduct invariance and of
\(\mathrm{Eq}\) by part 1 of \cref{thm:p02-w02-p02-005-bw-soundness}.
\Cref{thm:p02-w02-p02-005-bw-lemma-43} for \(\mathfrak M_C\) gives a truth
assignment satisfying \(T\cup T_{\mathrm{Henkin}}\cup\mathrm{Eq}\).

(iii)\(\Rightarrow\)(ii).  This follows Barwise's construction
\cite[ch.~A.1, \S4, pp.~31--33]{Barwise1977FirstOrder}; Steps 2, 3 and 5
supply the treatment of closed terms other than constants, which his text
passes over.  Let \(v\) satisfy
\(\Theta=T\cup T_{\mathrm{Henkin}}\cup\mathrm{Eq}\), and write
\(\Vdash\sigma\) for \(\bar v(\sigma)=\mathbf t\), \(\sigma\) a sentence of
\(L(C)\).  ``Closed term'' means closed term of \(L(C)\).

Step 1.  For \(\varphi(v)\) with \(\FV(\varphi)\subseteq\{v\}\):
(a) \(\Vdash\exists v\,\varphi(v)\) iff \(\Vdash\varphi(c_{\varphi(v)})\),
iff \(\Vdash\varphi(t)\) for some closed \(t\); (b)
\(\Vdash\forall v\,\varphi(v)\) iff \(\Vdash\varphi(c)\) for all \(c\in C\),
iff \(\Vdash\varphi(t)\) for all closed \(t\).  For (a): the type I axiom
gives the first ``only if''; a type IV axiom gives ``if'' from any closed
\(t\), in particular from \(c_{\varphi(v)}\).  For (b): type III gives
\(\Vdash\varphi(t)\) for every closed \(t\), in particular every
\(c\in C\).  If \(\Vdash\varphi(c)\) for every \(c\in C\) but not
\(\Vdash\forall v\,\varphi(v)\), the type II axiom
\(\varphi(c_{\neg\varphi(v)})\to\forall v\,\varphi(v)\) has value
\(\mathbf t\) with false consequent, so
\(\bar v(\varphi(c_{\neg\varphi(v)}))=\mathbf f\), a contradiction.

Step 2 (transfer from constants to closed terms).  Let \(\psi\) be a
formula with \(\FV(\psi)\subseteq\{v_1,\ldots,v_n\}\), the \(v_i\)
distinct.  If \(\Vdash\psi(c_1,\ldots,c_n)\) for all
\(c_1,\ldots,c_n\in C\), then \(\Vdash\psi(t_1,\ldots,t_n)\) for all closed
\(t_1,\ldots,t_n\).  Induction on \(n\); \(n=0\) is trivial.  For
\(n\geq1\) fix \(c_1\in C\); the formula \(\psi(c_1,v_2,\ldots,v_n)\)
satisfies the hypothesis for \(n-1\), so
\(\Vdash\psi(c_1,t_2,\ldots,t_n)\) for all closed \(t_2,\ldots,t_n\).  Now
fix closed \(t_2,\ldots,t_n\); then \(\chi(v_1)=\psi(v_1,t_2,\ldots,t_n)\)
has \(\FV(\chi)\subseteq\{v_1\}\) and \(\Vdash\chi(c_1)\) for all
\(c_1\in C\), so Step 1(b) gives \(\Vdash\chi(t_1)\) for every closed
\(t_1\).  Substituting closed terms never captures a variable.

Step 3 (equality).  For closed \(s,t\) put \(s\approx t\) iff
\(\Vdash(s=t)\).  Each of the following formulas has all its instances
with constants of \(C\) in \(\mathrm{Eq}\), since those instances are
equality axioms of \(L(C)\) without variables:
\((v=v)\); \((v=w)\to(w=v)\); \((v_1=v_2\land v_2=v_3)\to(v_1=v_3)\);
\((v_1=w_1\land\cdots\land v_n=w_n)\to(R(v_1\cdots v_n)\to R(w_1\cdots w_n))\);
\((v_1=w_1\land\cdots\land v_n=w_n)\to(f(v_1\cdots v_n)=f(w_1\cdots w_n))\),
the last with the term \(f(x_1\cdots x_n)\).  By Step 2 all their closed
instances have value \(\mathbf t\); reading them with the truth table,
\(\approx\) is an equivalence relation on closed terms, \(s_i\approx t_i\)
(\(i\leq n\)) implies \(\Vdash R(s_1\cdots s_n)\) iff
\(\Vdash R(t_1\cdots t_n)\) (use the axiom in both directions with
symmetry), and \(s_i\approx t_i\) implies
\(f(s_1\cdots s_n)\approx f(t_1\cdots t_n)\).  Moreover every closed \(t\)
has some \(c\in C\) with \(t\approx c\): for \(\varphi(v)=(t=v)\), the
type IV axiom \((t=t)\to\exists v\,(t=v)\) and \(\Vdash(t=t)\) give
\(\Vdash\exists v\,(t=v)\), and type I gives \(\Vdash(t=c_{\varphi(v)})\).

Step 4 (the structure).  Let \(M\) be the set of classes \([c]\) of
\(\approx\) restricted to \(C\) (nonempty, since \(c_{(v=v)}\in C\)).  Put
\(([c_1],\ldots,[c_n])\in R^{\mathfrak M}\) iff
\(\Vdash R(c_1\cdots c_n)\), well defined by Step 3.  Put
\(f^{\mathfrak M}([c_1],\ldots,[c_n])=[d]\) for any \(d\in C\) with
\(f(c_1\cdots c_n)\approx d\); such \(d\) exists and \([d]\) does not
depend on the choices, by Step 3.  For a constant symbol \(e\) of
\(L(C)\) put \(e^{\mathfrak M}=[d]\) for any \(d\in C\) with \(e\approx d\);
for \(e\in C\) this is \([e]\).  So \(M=\{c^{\mathfrak M}\mid c\in C\}\):
\(\mathfrak M\) is canonical.

Step 5 (closed terms).  If \(t\) is closed, \(d\in C\) and \(t\approx d\),
then \(t^{\mathfrak M}=[d]\).  Induction on \(t\).  For a constant \(e\),
\(e^{\mathfrak M}=[d']\) with \(e\approx d'\), and \(e\approx d\) gives
\(d\approx d'\).  For \(t=f(t_1\cdots t_n)\) choose \(d_i\in C\) with
\(t_i\approx d_i\); by induction \(t_i^{\mathfrak M}=[d_i]\), so
\(t^{\mathfrak M}=f^{\mathfrak M}([d_1],\ldots,[d_n])=[d']\) with
\(f(d_1\cdots d_n)\approx d'\); by Step 3
\(t\approx f(d_1\cdots d_n)\approx d'\), so \(d\approx d'\).

Step 6 (atomic sentences).  With \(s\approx d\), \(t\approx d'\),
\(d,d'\in C\): \(\mathfrak M\models(s=t)\) iff \([d]=[d']\) iff
\(d\approx d'\) iff \(s\approx t\).  With \(t_i\approx d_i\):
\(\mathfrak M\models R(t_1\cdots t_n)\) iff \(\Vdash R(d_1\cdots d_n)\) iff
\(\Vdash R(t_1\cdots t_n)\).

Step 7.  For every sentence \(\sigma\) of \(L(C)\),
\(\mathfrak M\models\sigma\) iff \(\Vdash\sigma\); induction on the number
of occurrences of \(\neg,\land,\lor,\to,\exists,\forall\), which is the
same for \(\varphi(c)\) as for \(\varphi(v)\).  Atomic sentences are Step 6
and connectives follow from the truth table.  Let
\(\sigma=\exists v\,\varphi(v)\).  If \(\Vdash\sigma\), Step 1(a) gives
\(\Vdash\varphi(c)\) for \(c=c_{\varphi(v)}\); by induction
\(\mathfrak M\models\varphi(c)\), so \(\mathfrak M\models\sigma\).  If
\(\mathfrak M\models\sigma\), some \([c]\) satisfies \(\varphi\); since
\([c]=c^{\mathfrak M}\), \cref{thm:fol-substitution} gives
\(\mathfrak M\models\varphi(c)\), so \(\Vdash\varphi(c)\) by induction and
\(\Vdash\sigma\) by Step 1(a).  Let \(\sigma=\forall v\,\varphi(v)\).
Since every element is some \(c^{\mathfrak M}\),
\(\mathfrak M\models\sigma\) iff \(\mathfrak M\models\varphi(c)\) for all
\(c\in C\), iff \(\Vdash\varphi(c)\) for all \(c\in C\) by induction, iff
\(\Vdash\sigma\) by Step 1(b).  In particular \(\mathfrak M\) is a
canonical model of \(T\).
\end{proof}
\end{annexitem}

\begin{annexitem}{annex:p02-w02-p02-005-bw-vacuous}{thm:p02-w02-p02-005-bw-vacuous}{The Main Lemma needs Henkin axioms for vacuous quantifiers}
\begin{proof}
No structure satisfies \(\neg(d=d)\) at any element, so \(T\) has no model.
Let \(\mathfrak M\) be the one-element \(L(C)\)-structure; it has a single
assignment, and \(L(C)\) has no relation symbols, so every atomic formula
is satisfied.  Define \(\models'\) by the clauses of satisfaction in
\(\mathfrak M\), except that \(\models'\exists v\,\varphi\) holds whenever
\(v\notin\FV(\varphi)\).  By induction on \(\varphi\), for every closed
term \(t\), \(\models'\varphi(t/v)\) iff \(\models'\varphi\): atomic
formulas are all satisfied; connectives are immediate; for
\(\forall u\,\psi\) and for \(\exists u\,\psi\) with \(u\neq v\), the
variable \(u\) is free in \(\psi(t/v)\) iff it is free in \(\psi\), because
\(t\) is closed, and the induction hypothesis applies to \(\psi\); for
\(u=v\) nothing is substituted.  Let \(w\) give each prime sentence
\(\pi\) the value \(\mathbf t\) iff \(\models'\pi\); since \(\models'\) is
truth-functional on connectives, \(\bar w(\sigma)=\mathbf t\) iff
\(\models'\sigma\) for every sentence.  Now \(\models'\exists x\,\neg(d=d)\)
since \(x\notin\FV(\neg(d=d))\), so \(w\) satisfies \(T\); it satisfies
\(\mathrm{Eq}\), all of whose members are built from atomic sentences.  For
\(\varphi(v)\) with exactly one free variable, \(\models'\exists v\,\varphi\)
iff \(\models'\varphi\) (one element), iff \(\models'\varphi(c)\) by the
substitution fact, and \(\models'\forall v\,\varphi\) iff
\(\models'\varphi(c)\) likewise, so the Henkin axioms of types I and II
hold.  Type III: \(\models'\forall v\,\varphi\) iff \(\models'\varphi\), iff
\(\models'\varphi(t)\).  Type IV: \(\models'\varphi(t)\) implies
\(\models'\varphi\), which implies \(\models'\exists v\,\varphi\) when
\(v\in\FV(\varphi)\); when \(v\notin\FV(\varphi)\) the consequent holds
outright.  So \(w\) satisfies \(T\cup T_{\mathrm{Henkin}}\cup\mathrm{Eq}\).
Under the reading of \cref{def:p02-w02-p02-005-bw-witnessing}, the type I
axiom \(\exists x\,\neg(d=d)\to\neg(d=d)\) and \((d=d)\in\mathrm{Eq}\) rule
this assignment out.
\end{proof}
\end{annexitem}

\begin{annexitem}{annex:p02-w02-p02-005-bw-soundness}{thm:p02-w02-p02-005-bw-soundness}{Soundness Lemma}
\begin{proof}
Part 1.  Tautologies: \cref{thm:p02-w02-p02-005-bw-lemma-43}.  Equality
axioms: under \(s\) each of \(u,w,u_i,w_i\) denotes an element; equality
is reflexive, symmetric and transitive; if \(u_i,w_i\) denote the same
elements, then so do \(t(u_1\cdots u_n)\) and \(t(w_1\cdots w_n)\)
(\cref{lem:term-substitution}), and \(R(u_1\cdots u_n)\),
\(R(w_1\cdots w_n)\) have the same value.  Quantifier axioms: if
\(\mathfrak M\models\forall v\,\varphi(v)[s]\), then
\(\mathfrak M\models\varphi[s(\tfrac{a}{v})]\) for
\(a=t^{\mathfrak M}(s)\), which is \(\mathfrak M\models\varphi(t)[s]\) by
\cref{thm:fol-substitution}; and \(\mathfrak M\models\varphi(t)[s]\) gives
\(\mathfrak M\models\varphi[s(\tfrac av)]\) for that \(a\), so
\(\mathfrak M\models\exists v\,\varphi(v)[s]\).

Part 2, Barwise's induction \cite[ch.~A.1, \S4, 4.10, p.~35]{Barwise1977FirstOrder}.
Let \(\psi_1,\ldots,\psi_n\) be a proof from \(T\).  Show by induction on
\(k\) that \(\mathfrak M\models\psi_k[s]\) for every assignment \(s\);
this is \(\mathfrak M\models\forall v_1\cdots\forall v_m\,\psi_k\) for the
free variables \(v_i\) of \(\psi_k\).  Axioms: part 1.  Members of \(T\):
sentences true in \(\mathfrak M\).  Modus ponens: immediate.  First
generalization rule, from \(\varphi\to\psi(v)\) with \(v\notin\FV(\varphi)\):
let \(\mathfrak M\models\varphi[s]\) and \(a\in M\).  The assignments
\(s\) and \(s(\tfrac av)\) agree on \(\FV(\varphi)\), so
\(\mathfrak M\models\varphi[s(\tfrac av)]\) (\cref{lem:fol-coincidence}) and
so \(\mathfrak M\models\psi(v)[s(\tfrac av)]\).  If \(y\) is \(v\), this is
\(\mathfrak M\models\psi(y)[s(\tfrac ay)]\).  Otherwise \(y\) is not free
in \(\psi(v)\) and is substitutable for \(v\), and
\cref{thm:fol-substitution} gives \(\mathfrak M\models\psi(y)[s(\tfrac ay)]\)
iff \(\mathfrak M\models\psi(v)[s(\tfrac ay)(\tfrac av)]\), iff
\(\mathfrak M\models\psi(v)[s(\tfrac av)]\) by \cref{lem:fol-coincidence}.
As \(a\) was arbitrary, \(\mathfrak M\models\forall y\,\psi(y)[s]\).  The
second generalization rule is dual: if
\(\mathfrak M\models\exists y\,\psi(y)[s]\), the same computation gives
\(a\) with \(\mathfrak M\models\psi(v)[s(\tfrac av)]\), hence
\(\mathfrak M\models\varphi[s(\tfrac av)]\) and
\(\mathfrak M\models\varphi[s]\).
\end{proof}
\end{annexitem}

\begin{annexitem}{annex:p02-w02-p02-005-bw-lemma-411}{thm:p02-w02-p02-005-bw-lemma-411}{Derived rules of H}
\begin{proof}
The concatenation of proofs from \(T\) is a proof from \(T\), and if
\(T\vdash\chi_1,\ldots,T\vdash\chi_m\) and
\(\chi_1\to\cdots\to\chi_m\to\chi\) is a tautology, then \(T\vdash\chi\) by
\(m\) applications of modus ponens; call this \emph{tautological
consequence}.  Each formula said to be a tautology below is checked by
assuming its consequent false.

(i) \([(\varphi\to\psi)\to((\neg\varphi\to\psi)\to\psi)]\) is a tautology:
with \(\psi\) false, the two antecedents require \(\varphi\) false and
\(\varphi\) true.  (ii) \([((\varphi\to\theta)\to\psi)\to(\neg\varphi\to\psi)]\)
and \([((\varphi\to\theta)\to\psi)\to(\theta\to\psi)]\) are tautologies:
\(\neg\varphi\) true, or \(\theta\) true, makes \(\varphi\to\theta\) true.

(iii) From the hypothesis, (ii) gives \(T\vdash\neg\exists y\,\varphi(y)\to\psi\)
and \(T\vdash\varphi(v)\to\psi\).  As \(v\) is not free in \(\psi\), the
second generalization rule gives \(T\vdash\exists y\,\varphi(y)\to\psi\),
and (i) gives \(T\vdash\psi\) \cite[ch.~A.1, \S4, p.~36]{Barwise1977FirstOrder}.

(iv) Barwise says only that it is similar, with the first rule.  By (ii),
\(T\vdash\neg\varphi(v)\to\psi\) and \(T\vdash\forall y\,\varphi(y)\to\psi\).
By tautological consequence \(T\vdash\neg\psi\to\varphi(v)\); as \(v\) is
not free in \(\neg\psi\), the first rule gives
\(T\vdash\neg\psi\to\forall y\,\varphi(y)\).  The tautology
\((\neg\psi\to B)\to((B\to\psi)\to\psi)\), with \(B=\forall y\,\varphi(y)\),
gives \(T\vdash\psi\).

(v) Put \(A=\exists y\,\neg\varphi(y)\), \(B=\forall y\,\varphi(y)\),
\(P=\varphi(v)\), \(I=(A\to\neg P)\), \(I\!I=(P\to B)\); the hypothesis is
\(T\vdash I\to(I\!I\to\psi)\).  Neither \(A\) nor \(B\) has \(v\) free.
The formulas \([I\to(I\!I\to\psi)]\to(\neg P\to\psi)\) and
\([I\to(I\!I\to\psi)]\to(P\to(\neg A\to(B\to\psi)))\) are tautologies:
\(\neg P\) true makes \(I\) and \(I\!I\) true, and \(P,\neg A,B\) true make
\(I\) and \(I\!I\) true.  So \(T\vdash\neg\varphi(v)\to\psi\) and
\(T\vdash\varphi(v)\to\chi\) with \(\chi=(\neg A\to(B\to\psi))\).  As \(v\)
is free in neither \(\psi\) nor \(\chi\), the second rule gives
\(T\vdash A\to\psi\) and \(T\vdash\exists y\,\varphi(y)\to\chi\).  The H
axiom \(\neg\varphi(v)\to\exists y\,\neg\varphi(y)\) (axiom (3) for the
formula \(\neg\varphi(y)\) and the term \(v\)) gives, by tautological
consequence, \(T\vdash\neg A\to\varphi(v)\), and the first rule gives
\(T\vdash\neg A\to B\).  The H axioms \(B\to\varphi(v)\) and
\(\varphi(v)\to\exists y\,\varphi(y)\) give
\(T\vdash B\to\exists y\,\varphi(y)\).  Finally
\((A\to\psi)\to((\neg A\to B)\to((B\to E)\to((E\to\chi)\to\psi)))\), with
\(E=\exists y\,\varphi(y)\), is a tautology: if \(A\) is true the first
antecedent gives \(\psi\); otherwise \(B\), \(E\), \(\chi\) and so \(\psi\)
are true in turn.  Hence \(T\vdash\psi\).
\end{proof}
\end{annexitem}

\begin{annexitem}{annex:p02-w02-p02-005-bw-completeness-h}{thm:p02-w02-p02-005-bw-completeness-h}{G\"odel's Completeness Theorem for H}
\begin{proof}
If \(T\vdash\varphi\), part 2 of \cref{thm:p02-w02-p02-005-bw-soundness}
with no free variables gives \(\mathfrak M\models\varphi\) for every model
\(\mathfrak M\) of \(T\).

Conversely suppose \(T\models\varphi\), following Barwise
\cite[ch.~A.1, \S4, pp.~36--37]{Barwise1977FirstOrder} with the pairing of
Henkin axioms added.  \(T\cup\{\neg\varphi\}\) has no model, so by the Main
Lemma \(T\cup\{\neg\varphi\}\cup T_{\mathrm{Henkin}}\cup\mathrm{Eq}\) is not
consistent in the sense of propositional logic, and by
\cref{thm:p02-w02-p02-005-bw-prop-compactness} there is a finite
\(S\subseteq T\cup T_{\mathrm{Henkin}}\cup\mathrm{Eq}\) with
\(S\cup\{\neg\varphi\}\) inconsistent.

Ordering.  Let \(\alpha_1,\ldots,\alpha_N\) be the members of \(S\) in
\(T\cup\mathrm{Eq}\) or of type III or IV.  The \emph{constant} of a Henkin
axiom is \(c_{\theta(v)}\) for the type I axiom of \(\theta(v)\) and
\(c_{\neg\theta(v)}\) for the type II axiom of \(\theta(v)\).  Two Henkin
axioms share a constant only if they are the type I axiom of
\(\neg\theta(v)\) and the type II axiom of \(\theta(v)\).  The rank of a
formula is the least \(n\) with the formula in \(L_n\).  If
\(c=c_{\theta(v)}\in C_{m+1}-C_m\), then \(\theta\in L_m-L_{m-1}\), and each
Henkin axiom with constant \(c\) lies in \(L_{m+1}\) but not in \(L_m\): it
has rank \(m+1\).  If a Henkin axiom with a constant \(d\neq c\) contains
\(c\), then \(c\) occurs in the formula \(\theta'\) of \(d\), so
\(\theta'\in L_k\) only for \(k\geq m+1\), and the axiom has rank at least
\(m+2\).  Group the Henkin axioms of \(S\) into \emph{blocks}, a block being
the set of Henkin axioms of \(S\) with a given constant (one or two
axioms).  Choose a block \(B_1\) of maximal rank, then a block \(B_2\) of
maximal rank among the rest, and so on, ending with \(B_m\).  By the rank
computation, the constant \(c_j\) of \(B_j\) occurs in no axiom of
\(B_{j+1},\ldots,B_m\), and it does not occur in the \(L\)-sentence
\(\varphi\).  List the Henkin axioms as \(\beta_1,\ldots,\beta_M\) block by
block, a two-element block in the order type I, type II.

The tautology.  Associating to the right,
\(\alpha_1\to\cdots\to\alpha_N\to\beta_1\to\cdots\to\beta_M\to\varphi\) is
a tautology, since an assignment making it false satisfies
\(S\cup\{\neg\varphi\}\).  Replace each witnessing constant \(c\) occurring
in it by a distinct variable \(x_c\) occurring nowhere in it, and write
\(\chi'\) for the result of this replacement on \(\chi\).  The replacement
sends prime formulas to prime formulas, commutes with the connectives, and
is undone by replacing each \(x_c\) by \(c\), so it is injective on the
prime constituents.  Given an assignment \(w\) to the primed prime
formulas, \(v(\pi)=w(\pi')\) defines an assignment with
\(\bar v(\chi)=\bar w(\chi')\) for every constituent \(\chi\); hence
\(\alpha_1'\to\cdots\to\alpha_N'\to\beta_1'\to\cdots\to\beta_M'\to\varphi\)
is a tautology, with \(\varphi'=\varphi\).

The \(\alpha_i'\) are provable.  If \(\alpha_i\in T\), it has no witnessing
constant and \(\alpha_i'=\alpha_i\in T\).  If \(\alpha_i\in\mathrm{Eq}\),
then \(\alpha_i'\) is an equality axiom of \(L\), the replaced constants
having become variables, except that in the last form the term may contain
variables \(x_{e_1},\ldots,x_{e_k}\) besides the \(u_i\); then
\(\alpha_i'\) follows tautologically from the equality axiom for the
\((n+k)\)-ary term with the extra pairs \(x_{e_j}=x_{e_j}\) and from the
axioms \((x_{e_j}=x_{e_j})\).  If \(\alpha_i\) has type III or IV, then
\(\alpha_i'\) is an axiom (3) of H: replacing constants by new variables
commutes with substitution, and no \(x_c\) is bound anywhere.  By \(N\)
applications of modus ponens, \(T\vdash\beta_1'\to\cdots\to\beta_M'\to\varphi\).

Removing the blocks.  Let \(\psi_j\) be the part of the chain after the
block \(B_j'\), so \(\psi_m=\varphi\).  The variable \(x_{c_j}\) does not
appear in \(\psi_j\).  If \(B_j\) is a single type I axiom
\((\exists v\,\theta(v))\to\theta(c_j)\), then \(B_j'\) is
\((\exists v\,\theta'(v))\to\theta'(x_{c_j})\), where \(v\) is not free in
\(\theta'(x_{c_j})\) and is substitutable for \(x_{c_j}\) in it (the
\(x_{c_j}\) occupy the places where \(v\) was free); 4.11(iii) passes from
\(T\vdash B_j'\to\psi_j\) to \(T\vdash\psi_j\).  A single type II axiom is
removed by 4.11(iv) in the same way, and a two-element block by clause (v)
of \cref{thm:p02-w02-p02-005-bw-lemma-411}.  After \(m\) steps,
\(T\vdash\varphi\).
\end{proof}
\end{annexitem}

\begin{annexitem}{annex:p02-w02-p02-005-bw-g-prop-completeness}{thm:p02-w02-p02-005-bw-g-prop-completeness}{Completeness of the propositional fragment of G}
\begin{proof}
Call \(\Gamma\vdash\Delta\) \emph{valid} if (i) holds for it, and say that
an assignment \(v\) \emph{refutes} it if \(v\) makes every member of
\(\Gamma\) true and every member of \(\Delta\) false.

(ii)\(\Rightarrow\)(i), by induction on the derivation.  No assignment
refutes an axiom \(\Gamma,\varphi\vdash\Delta,\varphi\).  For each rule, an
assignment refuting the conclusion refutes some premise: for
\((\land\vdash)\), \(\varphi\land\psi\) true makes \(\varphi\) and \(\psi\)
true; for \((\vdash\land)\), \(\varphi\land\psi\) false makes \(\varphi\)
or \(\psi\) false; for \((\lor\vdash)\), \(\varphi\lor\psi\) true makes
\(\varphi\) or \(\psi\) true; for \((\vdash\lor)\), \(\varphi\lor\psi\)
false makes both false; for \((\neg\vdash)\) and \((\vdash\neg)\),
\(\neg\varphi\) true (false) makes \(\varphi\) false (true); for
\((\to\vdash)\), \(\varphi\to\psi\) true makes \(\varphi\) false or
\(\psi\) true; for \((\vdash\to)\), \(\varphi\to\psi\) false makes
\(\varphi\) true and \(\psi\) false.  So validity passes from premises to
conclusion.

(i)\(\Rightarrow\)(ii), Barwise's backward construction
\cite[ch.~A.1, \S4, p.~38]{Barwise1977FirstOrder}.  Starting from
\(\Gamma\vdash\Delta\), repeatedly choose a sequent at a leaf that contains
a non-prime formula, choose such a formula of maximal length, and write
above the sequent the premises of the rule that introduces that formula on
its side, with the formula removed from the side set: for instance, for
\((\psi\to\theta)\in\Gamma\) and \(\Gamma_0=\Gamma-\{(\psi\to\theta)\}\),
the premises \(\Gamma_0\vdash\Delta,\psi\) and \(\Gamma_0,\theta\vdash\Delta\),
from which \((\to\vdash)\) yields \(\Gamma_0,(\psi\to\theta)\vdash\Delta\).
Each step lowers the total number of connective occurrences in the leaf, so
the process stops, with leaves containing only prime formulas.  If every
leaf has a prime formula on both sides, it is an axiom and the tree is a
derivation.  Otherwise some leaf \(\Gamma'\vdash\Delta'\) has
\(\Gamma'\cap\Delta'=\emptyset\); let \(v(p)=\mathbf t\) for
\(p\in\Gamma'\), \(v(p)=\mathbf f\) for \(p\in\Delta'\), and arbitrary
elsewhere.  Then \(v\) refutes the leaf.  For each step of the
construction, an assignment refuting one of the written premises refutes
the sequent below, by the converse computations: for the example step,
\(\psi\) false, or \(\theta\) true, makes \(\psi\to\theta\) true.  Passing
down the tree, \(v\) refutes \(\Gamma\vdash\Delta\), contradicting (i).
\end{proof}
\end{annexitem}

\begin{annexitem}{annex:p02-w02-p02-005-bw-completeness-g}{thm:p02-w02-p02-005-bw-completeness-g}{Completeness Theorem for G\textsuperscript{+}}
\begin{proof}
\emph{Soundness} (\(\Rightarrow\)), the analogue of 4.10 that Barwise
invokes \cite[ch.~A.1, \S4, p.~39]{Barwise1977FirstOrder}.  Call
\(\Gamma\vdash\Delta\) \emph{valid} if every assignment in every structure
that satisfies all of \(\Gamma\) satisfies some member of \(\Delta\).
Axioms are valid, including \(\Gamma\vdash\Delta,(t=t)\).  The
propositional rules and cut preserve validity by the computations in the
proof of \cref{thm:p02-w02-p02-005-bw-g-prop-completeness}, applied to the
truth assignment of \cref{thm:p02-w02-p02-005-bw-lemma-43}.  Equality rule:
if \(s\) satisfies \(E\), then \(t_1^{\mathfrak M}(s)=t_2^{\mathfrak M}(s)\),
so by \cref{thm:fol-substitution} \(s\) satisfies \(\varphi(t_1)\) iff it
satisfies \(\varphi(t_2)\), and likewise for \(\psi\).  \((\forall\vdash)\)
and \((\vdash\exists)\): \cref{thm:fol-substitution} as in part 1 of
\cref{thm:p02-w02-p02-005-bw-soundness}.  \((\vdash\forall)\): let \(s\)
satisfy \(\Gamma\) and no member of \(\Delta\).  For \(a\in M\), the
assignment \(s(\tfrac av)\) agrees with \(s\) on the free variables of
\(\Gamma\cup\Delta\), so it satisfies \(\Gamma\) and no member of \(\Delta\)
(\cref{lem:fol-coincidence}), hence satisfies \(\varphi(v)\) by validity of
the premise; the computation in the proof of
\cref{thm:p02-w02-p02-005-bw-soundness} then gives \(\forall y\,\varphi(y)\)
at \(s\).  \((\exists\vdash)\) is dual.  So derivable sequents are valid,
and if \(\Gamma\vdash\{\varphi\}\) is derivable with \(\Gamma\subseteq T\)
finite, every model of \(T\) satisfies \(\varphi\).

\emph{Completeness} (\(\Leftarrow\)).  Suppose \(T\models\varphi\).  As in
the proof of \cref{thm:p02-w02-p02-005-bw-completeness-h}, there is a finite
\(S\subseteq T\cup T_{\mathrm{Henkin}}\cup\mathrm{Eq}\) with
\(S\cup\{\neg\varphi\}\) propositionally inconsistent; list it as
\(\alpha_1,\ldots,\alpha_N\) and blocks \(B_1,\ldots,B_m\) of Henkin axioms,
and form the primed formulas.  Every truth assignment satisfying
\(\Gamma=\{\alpha_1',\ldots,\alpha_N'\}\cup B_1'\cup\cdots\cup B_m'\) makes
\(\varphi\) true, so by \cref{thm:p02-w02-p02-005-bw-g-prop-completeness},
with the atomic and quantified formulas as prime formulas,
\(\Gamma\vdash\varphi\) has a derivation using only the propositional
axioms and rules.  Let \(\Gamma_0\) be the set of \(\alpha_i'\) that belong
to \(T\).  The \emph{eigenvariables} of a derivation are the variables
\(v\) of its applications of \((\vdash\forall)\) and \((\exists\vdash)\).
Barwise's claims (1)--(5) \cite[ch.~A.1, \S4, pp.~39--40]{Barwise1977FirstOrder}
and a sixth are proved as follows.

(1) If \(\alpha\) is an equality axiom or a quantifier axiom, possibly
with some constants replaced by variables, then \(\emptyset\vdash\alpha\)
has a derivation without eigenvariables.  Quantifier axioms:
\(\forall v\,\theta(v)\vdash\theta(t)\) follows by \((\forall\vdash)\) from
the axiom \(\theta(t)\vdash\theta(t)\), and \((\vdash\to)\) gives
\(\emptyset\vdash\forall v\,\theta(v)\to\theta(t)\); type IV uses
\((\vdash\exists)\) in the same way.  For equality first derive the rule
(E\('\)): from \(\Gamma\vdash\Delta,\psi(t_1)\), with a derivation without
eigenvariables, infer \(\Gamma,E\vdash\Delta,\psi(t_2)\) for \(E\) either
\((t_1=t_2)\) or \((t_2=t_1)\).  Add \((t_1=t_1)\) to every left side of the
given derivation (an instance of (2) below), apply the equality rule with
\(\varphi(x)=(x=x)\) to get \(\Gamma,E,(t_2=t_2)\vdash\Delta,\psi(t_2)\), and
cut \((t_2=t_2)\) against the equality axiom
\(\Gamma,E\vdash\Delta,\psi(t_2),(t_2=t_2)\).  Then \(\vdash(u=u)\) is an
equality axiom; \((u=w)\vdash(w=u)\) comes by (E\('\)) from
\(\vdash(u=u)\) with \(\psi(x)=(x=u)\), \(t_1=u\), \(t_2=w\);
\((u_1=u_2),(u_2=u_3)\vdash(u_1=u_3)\) comes by (E\('\)) from the axiom
\((u_1=u_2)\vdash(u_1=u_2)\) with \(\psi(x)=(u_1=x)\);
\(R(u_1\cdots u_n)\) and \(E_i=(u_i=w_i)\), \(i\leq n\), yield
\(R(w_1\cdots w_n)\) by \(n\) applications of (E\('\)) to the axiom
\(R(u_1\cdots u_n)\vdash R(u_1\cdots u_n)\), the \(i\)-th with
\(\psi(x)=R(w_1\cdots w_{i-1}xu_{i+1}\cdots u_n)\); and the term axiom comes
in the same way from \(\vdash t(u_1\cdots u_n)=t(u_1\cdots u_n)\) with
\(\psi(x)=(t(u_1\cdots u_n)=t(w_1\cdots w_{i-1}xu_{i+1}\cdots u_n))\), any
further variables of \(t\) being left in place.  In each case
\((\land\vdash)\) collects the hypotheses into the right-associated
conjunction and \((\vdash\to)\) forms the implication.

(2) Weakening.  If \(D\) derives \(\Gamma\vdash\Delta\) and no
eigenvariable of \(D\) is free in \(\Gamma'\cup\Delta'\), then adding
\(\Gamma'\) to every left side and \(\Delta'\) to every right side of \(D\)
gives a derivation of \(\Gamma\cup\Gamma'\vdash\Delta\cup\Delta'\) with
the same eigenvariables: each axiom and each rule instance stays one,
since the side sets are enlarged uniformly and the eigenvariable
conditions still hold.  Barwise states (2) without the condition; it is
needed, since \(\varphi(v)\vdash\forall y\,\varphi(y)\) may not be
weakened by adding a formula with \(v\) free on its left.

Applying (1) and (2), every \(\alpha_i'\notin\Gamma_0\) has a derivation of
\(\Gamma''\vdash\varphi,\alpha_i'\) without eigenvariables for any
\(\Gamma''\), and \(N-K\) cuts, \(K=\abs{\Gamma_0}\), turn the derivation of
\(\Gamma\vdash\varphi\) into one of
\(\Gamma_0\cup B_1'\cup\cdots\cup B_m'\vdash\varphi\) without
eigenvariables.  (Barwise writes \(N-J\) and \(J<i\leq N\) at this point;
the count is over the \(\alpha_i'\) outside \(\Gamma_0\).)

In (3)--(6) let \(D\) be the given derivation, and assume its eigenvariables
are not free in the formulas named in the claim; this holds below because
the eigenvariables are \(x_{c_1},\ldots,x_{c_{j-1}}\) and fresh variables,
none of which occurs in \(\Gamma_0\), in \(B_j',\ldots,B_m'\) or in
\(\varphi\).

(3) If \(\Gamma,(\varphi\to\theta)\vdash\Delta\), then
\(\Gamma\vdash\Delta,\varphi\) and \(\Gamma,\theta\vdash\Delta\).  The first
is Barwise's derivation: weaken to
\(\Gamma,(\varphi\to\theta)\vdash\Delta,\varphi\), take the axiom
\(\Gamma,\varphi\vdash\Delta,\theta,\varphi\), apply \((\vdash\to)\) to get
\(\Gamma\vdash\Delta,(\varphi\to\theta),\varphi\), and cut.  The second:
weaken to \(\Gamma,\theta,(\varphi\to\theta)\vdash\Delta\); from the axiom
\(\Gamma,\theta,\varphi\vdash\Delta,\theta\), \((\vdash\to)\) gives
\(\Gamma,\theta\vdash\Delta,(\varphi\to\theta)\); cut.

(4) If \(\Gamma,(\exists y\,\varphi(y)\to\varphi(v))\vdash\Delta\) and \(v\)
is not free in \(\Gamma\cup\Delta\), then \(\Gamma\vdash\Delta\): (3) gives
\(\Gamma\vdash\Delta,\exists y\,\varphi(y)\) and
\(\Gamma,\varphi(v)\vdash\Delta\); \((\exists\vdash)\) gives
\(\Gamma,\exists y\,\varphi(y)\vdash\Delta\); cut.

(5) If \(\Gamma,(\varphi(v)\to\forall y\,\varphi(y))\vdash\Delta\) and \(v\)
is not free in \(\Gamma\cup\Delta\), then \(\Gamma\vdash\Delta\): (3) gives
\(\Gamma\vdash\Delta,\varphi(v)\) and \(\Gamma,\forall y\,\varphi(y)\vdash\Delta\);
\((\vdash\forall)\) gives \(\Gamma\vdash\Delta,\forall y\,\varphi(y)\); cut.

(6) Let \(A=\exists y\,\neg\varphi(y)\), \(B=\forall y\,\varphi(y)\),
\(I=(A\to\neg\varphi(v))\), \(I\!I=(\varphi(v)\to B)\).  If
\(\Gamma,I,I\!I\vdash\Delta\) and \(v\) is not free in \(\Gamma\cup\Delta\),
then \(\Gamma\vdash\Delta\).  First, \(\Gamma\vdash\Delta,A\).  Let \(w\) be
a variable occurring nowhere so far.  From the axiom
\(\Gamma,\varphi(v),\varphi(w)\vdash\Delta,\varphi(w)\), rule
\((\vdash\neg)\) gives \(\Gamma,\varphi(v)\vdash\Delta,\varphi(w),\neg\varphi(w)\),
\((\vdash\exists)\) gives \(\Gamma,\varphi(v)\vdash\Delta,\varphi(w),A\), and
\((\vdash\forall)\), with \(w\) not free in \(\Gamma,\varphi(v),\Delta,A\),
gives \(\Gamma,\varphi(v)\vdash\Delta,A,B\); \((\vdash\to)\) gives
\(\Gamma\vdash\Delta,A,I\!I\).  From the axiom
\(\Gamma,A\vdash\Delta,A,\neg\varphi(v)\), \((\vdash\to)\) gives
\(\Gamma\vdash\Delta,A,I\).  Weaken \(D\) to \(\Gamma,I,I\!I\vdash\Delta,A\)
and \(\Gamma\vdash\Delta,A,I\!I\) to \(\Gamma,I\vdash\Delta,A,I\!I\), cut
\(I\!I\), then cut \(I\).  Second, \(\Gamma,A\vdash\Delta\).  From the axiom
\(\Gamma,\neg\varphi(v),A\vdash\Delta,\neg\varphi(v)\), \((\vdash\to)\)
gives \(\Gamma,\neg\varphi(v)\vdash\Delta,I\).  From the axiom
\(\Gamma,\varphi(v)\vdash\Delta,B,\varphi(v)\), \((\neg\vdash)\) gives
\(\Gamma,\neg\varphi(v),\varphi(v)\vdash\Delta,B\) and \((\vdash\to)\)
gives \(\Gamma,\neg\varphi(v)\vdash\Delta,I\!I\).  Weaken \(D\) to
\(\Gamma,\neg\varphi(v),I,I\!I\vdash\Delta\) and cut \(I\!I\), then \(I\),
to get \(\Gamma,\neg\varphi(v)\vdash\Delta\); \((\exists\vdash)\), with
\(v\) not free in \(\Gamma\cup\Delta\), gives \(\Gamma,A\vdash\Delta\).
Cutting \(A\) gives \(\Gamma\vdash\Delta\).

Now remove the blocks in the order \(B_1,\ldots,B_m\), as in the proof of
\cref{thm:p02-w02-p02-005-bw-completeness-h}: at stage \(j\) take
\(\Gamma=\Gamma_0\cup B_{j+1}'\cup\cdots\cup B_m'\), \(\Delta=\{\varphi\}\)
and \(v=x_{c_j}\), which is not free in \(\Gamma\cup\Delta\); a type I
block is removed by (4), a type II block by (5) and a pair by (6).  The
result is a derivation of \(\Gamma_0\vdash\varphi\) with
\(\Gamma_0\subseteq T\) finite.
\end{proof}
\end{annexitem}

\begin{annexitem}{annex:p02-w02-p02-005-bw-exercise}{thm:p02-w02-p02-005-bw-exercise}{The exercise is provable in G and in H}
\begin{proof}
Write \(\theta(y)=(\varphi(y)\to\forall x\,\varphi(x))\) and
\(E=\exists y\,\theta(y)\), a sentence; let \(z\) be a variable not
occurring in \(\varphi(x)\) and distinct from \(y\).  In \(\mathbf G\),
reading each line from the one above:
\[
\begin{array}{ll}
\varphi(y),\varphi(z)\vdash E,\varphi(z),\forall x\,\varphi(x) & \text{axiom}\\
\varphi(y)\vdash E,\varphi(z),\theta(z) & (\vdash\to)\\
\varphi(y)\vdash E,\varphi(z) & (\vdash\exists),\ t=z\\
\varphi(y)\vdash E,\forall x\,\varphi(x) & (\vdash\forall),\ z\text{ not free in }\varphi(y),E\\
\vdash E,\theta(y) & (\vdash\to)\\
\vdash E & (\vdash\exists),\ t=y
\end{array}
\]
The two applications of \((\vdash\exists)\) have \(E\) already on the
right, which the set notation absorbs.  This follows Barwise's informal
proof: the case split on whether some \(y\) has \(\neg\varphi(y)\) becomes
the choice of the witness \(z\), produced by \((\vdash\forall)\).

In H: \(\theta(y)\to E\) is an axiom (3).  The tautologies
\(\neg\varphi(y)\to\theta(y)\) and \(\forall x\,\varphi(x)\to\theta(y)\)
give \(\vdash\neg\varphi(y)\to E\) and \(\vdash\forall x\,\varphi(x)\to E\)
by tautological consequence.  The second generalization rule, \(y\) not
free in \(E\), gives \(\vdash\exists y\,\neg\varphi(y)\to E\).  The axiom
\(\neg\varphi(x)\to\exists y\,\neg\varphi(y)\) gives
\(\vdash\neg\exists y\,\neg\varphi(y)\to\varphi(x)\), and the first rule,
\(x\) not free in the antecedent, gives
\(\vdash\neg\exists y\,\neg\varphi(y)\to\forall x\,\varphi(x)\).  The
tautology \((A\to E)\to((B\to E)\to((\neg A\to B)\to E))\) with
\(A=\exists y\,\neg\varphi(y)\), \(B=\forall x\,\varphi(x)\) gives
\(\vdash E\).
\end{proof}
\end{annexitem}

\subsection{Proofs for Barwise's \S5}
\label{sec:annex-p02-w02-p02-005-bw-s5}

\begin{annexitem}{annex:p02-w02-p02-005-bw-n-logic}{thm:p02-w02-p02-005-bw-n-logic}{Compactness fails for \(\mathfrak N\)-logic}
\begin{proof}
Let \(\kappa=\abs N\) and \(\lambda=\kappa^+\).  Use a function symbol
\(F\) from the first sort to the second and first-sort constants \(d_i\),
\(i<\lambda\).  Let \(T\) consist of \(d_i\neq d_j\) for \(i<j<\lambda\)
and \(\forall x\,\forall x'\,(F(x)=F(x')\to x=x')\).  A finite
\(T_0\subseteq T\) names finitely many \(d_i\); interpret the first sort as
\(N\), \(F\) as the identity, those \(d_i\) as distinct elements (\(N\) is
infinite) and the others arbitrarily; this is a model of \(T_0\) whose
second sort is \(\mathfrak N\).  A model of \(T\) would have at least
\(\lambda\) first-sort elements, mapped injectively by \(F\) into \(N\) of
size \(\kappa<\lambda\), which is impossible.  The source asserts the
failure without proof \cite[ch.~A.1, \S5, pp.~42--43]{Barwise1977FirstOrder}.
\end{proof}
\end{annexitem}

\begin{annexitem}{annex:p02-w02-p02-005-bw-weak-so-ls}{thm:p02-w02-p02-005-bw-weak-so-ls}{L\"owenheim--Skolem and compactness for weak second-order logic}
\begin{proof}
For \(b\in\mathrm{HF}(M)\) let \(\mathrm{supp}(b)=(b\cap M)\cup
\bigcup\{\mathrm{supp}(c)\mid c\in b\cap\mathrm{HF}(M)\}\), defined by
recursion on the least \(n\) with \(b\in\mathrm{HF}_n(M)\); it is finite.
For \(A\subseteq M\): \(\mathrm{HF}_n(A)\subseteq\mathrm{HF}_n(M)\) by
induction on \(n\), and if \(\mathrm{supp}(b)\subseteq A\), then
\(b\in\mathrm{HF}(A)\), by induction on the least \(n\) with
\(b\in\mathrm{HF}_n(M)\): the members of \(b\) in \(M\) lie in \(A\), and
the others have support inside \(\mathrm{supp}(b)\), so lie in
\(\mathrm{HF}_{n-1}(A)\).  If \(A\) is countable, so is each
\(\mathrm{HF}_n(A)\) and hence \(\mathrm{HF}(A)\).

Part 1.  The set of \(L^{*}\)-formulas is countable.  Fix a choice
function.  Let \(A_0\) be the closure of \(\{m_0\}\), some \(m_0\in M\),
together with the constants of \(L\), under the functions of
\(\mathfrak M\).  Given countable \(A_k\), for each formula
\(\exists x\,\psi\) or \(\exists a\,\psi\) and each assignment of elements
of \(A_k\) and of \(\mathrm{HF}(A_k)\) to its finitely many free variables
at which it is true in \(\mathrm{HF}(\mathfrak M)\), choose a witness: an
element of \(M\), or a set \(b\in\mathrm{HF}(M)\), in which case take
\(\mathrm{supp}(b)\).  Let \(A_{k+1}\) be the closure of \(A_k\) and these
countably many finite sets under the functions of \(\mathfrak M\), and
\(A=\bigcup_kA_k\), a countable set closed under the functions and
containing the constants; let \(\mathfrak A\) be the substructure on
\(A\).  Every \(b\in\mathrm{HF}(A)\) has finite support, contained in some
\(A_k\), so \(b\in\mathrm{HF}(A_k)\).  By induction on \(\psi\), for every
assignment \(s\) into \(A\) and \(\mathrm{HF}(A)\),
\(\mathrm{HF}(\mathfrak A)\models\psi[s]\) iff
\(\mathrm{HF}(\mathfrak M)\models\psi[s]\).  Atomic formulas of \(L\) have
the same value in the substructure; \(x\in a\), \(b\in a\) and equations
between sets have the same value because the sets are the same objects
and membership is restricted; connectives are immediate.  For
\(\exists x\,\psi\) and \(\exists a\,\psi\), a witness in
\(\mathrm{HF}(\mathfrak A)\) is one in \(\mathrm{HF}(\mathfrak M)\);
conversely the finitely many values of \(s\) on the free variables lie in
some \(A_k\cup\mathrm{HF}(A_k)\), so a witness was chosen in \(A_{k+1}\),
or in \(\mathrm{HF}(A_{k+1})\) by the support property, and the induction
hypothesis applies.  The universal quantifiers are the duals.  So
\(\mathrm{HF}(\mathfrak A)\) and \(\mathrm{HF}(\mathfrak M)\) satisfy the
same sentences, in particular those of \(T\).

Part 2.  Let \(\delta_n\) be the first-order sentence
\(\exists x_1\cdots\exists x_n\bigwedge_{i<j\leq n}x_i\neq x_j\), and
\(T=\{\exists a\,\forall x\,(x\in a)\}\cup\{\delta_n\mid n\geq2\}\).
\(\mathrm{HF}(\mathfrak M)\models\exists a\,\forall x\,(x\in a)\) iff \(M\)
is finite: if \(M\) is finite, then \(M\in\mathrm{HF}_1(M)\); if \(M\subseteq a\)
for a finite \(a\), then \(M\) is finite.  A finite subset of \(T\) mentions
\(\delta_n\) only for \(n\leq N\) and is true in \(\mathrm{HF}\) of an
\(N\)-element set; \(T\) would need a finite \(M\) with at least \(n\)
elements for every \(n\).  Barwise's own indication, at p.~10, is the proof
of his 2.1 \cite[ch.~A.1, \S2, p.~10]{Barwise1977FirstOrder}; the argument
above uses finiteness directly.
\end{proof}
\end{annexitem}

\begin{annexitem}{annex:p02-w02-p02-005-bw-countable-field}{thm:p02-w02-p02-005-bw-countable-field}{A countable field with the weak second-order properties of \(\mathbb R\)}
\begin{proof}
Apply part 1 of \cref{thm:p02-w02-p02-005-bw-weak-so-ls} to \(\mathbb R\)
and the set of \(L^{*}\)-sentences true in \(\mathrm{HF}(\mathbb R)\): it
gives a countable substructure \(\mathfrak A\) with
\(\mathrm{HF}(\mathfrak A)\) satisfying the same \(L^{*}\)-sentences.  A
first-order sentence of \(L\) is an \(L^{*}\)-sentence true in
\(\mathrm{HF}(\mathfrak M)\) iff true in \(\mathfrak M\), so \(\mathfrak A\)
satisfies the ordered field axioms.  Let \(\mathrm{Arch}^{*}\) be
\[
\forall x\,\exists a\,[0\in a\land\forall y\,(y\in a\to(y=0\lor
\exists z\,(z\in a\land z+1=y)))\land\exists y\,(y\in a\land(x<y\lor x=y))].
\]
In an ordered field \(\mathfrak F\), \(\mathrm{HF}(\mathfrak F)\models\mathrm{Arch}^{*}\)
iff \(\mathfrak F\) is Archimedean: if \(a\) is a finite set as described
and \(y\in a\), then \(y,y-1,y-2,\ldots\) stay in \(a\) while nonzero and
strictly decrease, so, \(a\) being finite, some \(y-n\cdot1\) is \(0\) and
\(y=n1\); conversely \(\{0,1,\ldots,n1\}\) witnesses \(x\leq n1\).  The
real field is Archimedean, so \(\mathrm{Arch}^{*}\) holds in
\(\mathrm{HF}(\mathbb R)\) and therefore in \(\mathrm{HF}(\mathfrak A)\).
\end{proof}
\end{annexitem}

\begin{annexitem}{annex:p02-w02-p02-005-bw-extended-ls}{thm:p02-w02-p02-005-bw-extended-ls}{L\"owenheim--Skolem and compactness for \(L_{\omega_1\omega}\), \(L(Q_0)\) and \(L(Q)\)}
\begin{proof}
Part 1.  Each formula of \(L_{\omega_1\omega}\) has countably many
subformulas (induction on formulas, using countable choice), so the set
\(F\) of subformulas of members of a countable \(T\) is countable.  Every
path through the formation tree of a sentence passes finitely many
quantifiers, so every member of \(F\) has finitely many free variables.
Let \(\mathfrak M\models T\).  Build countable sets
\(A_0\subseteq A_1\subseteq\cdots\subseteq M\) as in the proof of
\cref{thm:p02-w02-p02-005-bw-weak-so-ls}: \(A_0\) is the closure of one
element and the constants under the functions; \(A_{k+1}\) adds, for each
\(\exists x\,\psi\in F\) true at an assignment into \(A_k\) a witness, and
for each \(\forall x\,\psi\in F\) false at such an assignment a
counterexample, and closes under the functions.  Let \(\mathfrak A\) be the
substructure on \(A=\bigcup_kA_k\).  By induction along the subformula
relation, each \(\psi\in F\) has the same value in \(\mathfrak A\) and
\(\mathfrak M\) at every assignment into \(A\): atomic formulas and
connectives, including \(\bigwedge\) and \(\bigvee\), are immediate, and
the quantifier cases use the chosen witnesses and counterexamples.  So
\(\mathfrak A\) is a countable model of \(T\).

Part 2.  Translate \(Qx\,\varphi(x)\) into
\(\bigwedge_{n\geq1}\exists x_1\cdots\exists x_n\,(\bigwedge_{i<j\leq n}x_i\neq x_j
\land\bigwedge_{i\leq n}\varphi(x_i))\), with new variables, and leave the
other clauses unchanged.  By induction on formulas each formula of
\(L(Q_0)\) and its translation are satisfied by the same assignments in
every structure, since a set is infinite iff it has \(n\) distinct members
for every \(n\).  Apply part 1 to the translated set.

Part 3.  With \(\delta_n\) as in the proof of \cref{thm:p02-w02-p02-005-bw-weak-so-ls},
the sets \(\{\bigvee_{n\geq1}\exists x_1\cdots\exists x_n\,\forall y\,
(y=x_1\lor\cdots\lor y=x_n)\}\cup\{\delta_n\mid n\geq2\}\) of
\(L_{\omega_1\omega}\) and \(\{\neg Qx\,(x=x)\}\cup\{\delta_n\mid n\geq2\}\)
of \(L(Q_0)\) have finite models for every finite subset, but a model of
either would be finite and have at least \(n\) elements for each \(n\).

Part 4.  \(Qx\,(x=x)\) holds in a structure iff its domain is uncountable;
any uncountable structure is a model, and no countable one.
\end{proof}
\end{annexitem}

\begin{annexitem}{annex:p02-w02-p02-005-bw-keisler-axioms}{thm:p02-w02-p02-005-bw-keisler-axioms}{Keisler's axioms are valid and do not give ``many''}
\begin{proof}
Write the axioms as \(\forall y\,\neg Qx\,(x=y)\);
\(\forall x\,(\varphi(x)\to\psi(x))\to(Qx\,\varphi(x)\to Qx\,\psi(x))\);
\(Qx\,(\varphi(x)\lor\psi(x))\to(Qx\,\varphi(x)\lor Qx\,\psi(x))\); and
\((\neg Qx\,\exists y\,\varphi(x,y)\land\forall x\,\neg Qy\,\varphi(x,y))\to
\neg Qy\,\exists x\,\varphi(x,y)\), the formulas possibly having further
free variables.  Fix a structure and an assignment, and write
\(X_\chi\) for the set of elements satisfying \(\chi\) in the displayed
variable.  (1) \(X_{x=y}\) is a singleton.  (2) If
\(X_\varphi\subseteq X_\psi\) and \(X_\varphi\) is uncountable, so is
\(X_\psi\).  (3) \(X_{\varphi\lor\psi}=X_\varphi\cup X_\psi\), and a union
of two countable sets is countable.  (4)
\(\{y\mid\exists x\,\varphi(x,y)\}=\bigcup_{x\in D}\{y\mid\varphi(x,y)\}\)
with \(D=\{x\mid\exists y\,\varphi(x,y)\}\) countable and each fibre
countable; a countable union of countable sets is countable, by countable
choice.  So every instance holds everywhere, in particular in every
countable structure, where \(Qx\,(x=x)\) is false; hence \(Qx\,(x=x)\) is
not a consequence of the instances \cite[ch.~A.1, \S5, p.~45]{Barwise1977FirstOrder}.
\end{proof}
\end{annexitem}
